The university system is a disaster in just about every way possible.
However the \emph{potential} of this system is too great to ignore. The
potential of all those excited young people ready to change the world
who also happen to be on some of the nicest land in the world, in
institutions which have a fair amount of independence can truly change
the world if they change how they are run.

Much of this chapter should not be written by me but by the people who
are scholars of the history of student insurrections. Student
insurrection has been a critical component of both university life and
revolutionary practice for many centuries. There are whole movements
dedicated to this kind of work, some of which are documented in the
Undercommons Project at \url{http://undercommoning.org}. If you are on a
university campus, study the links contained there and study university
history and the history of student revolt globally. Study this history
and learn how to apply it to your own story and your own campus.

\subsection{Universities: Visions of
Utopia}\label{universities-visions-of-utopia}

If the university campus lived up to its potential it could be a true
paradise: essentially a giant garden filled with buildings for studying
and creating knowledge. Amazing! They are usually situated on some of
the best land to be found anywhere, have great access to everything
needed in life, and have dense urban style housing in a pastoral
environment which allows for a simple, car free life.

University campuses are often physically spectacular. They often have
some of the greatest examples of art and architecture of available both
used in their construction and lovingly maintained for in some cases
hundreds of years. It is typical for them to have wooded areas owned by
the university, as well as often running water and in many cases access
to very large bodies of water. University campuses are often more self
contained than most institutions, creating their own power and managing
their own utilities, and having a fair amount of autonomy from local
government.

\subsection{Undergraduate Education: Broken
Promises}\label{undergraduate-education-broken-promises}

The college education has become a key element in the great lie of the
American middle class dream. It is also a major factor in the older
generation destroying the opportunities they had for the younger
generation.

College has become just another capitalist cartel, exploiting the hopes
of young people in exchange for a life of debt servitude. Young people
are still told by their elders that a college education is needed to
enter the middle class, which is supposed to be a good thing. They're
told that the price is justified by long term earnings. Then they're
sent off to live in a party resort for 5 years where they shuffle
through a series of pointless and irrelevant classes, and wind up with a
bill of many tens of thousands of dollars(this is primarily an American
problem, but the neoliberals will bring it to your country soon enough
if they're not crushed at the global level).

And what do we learn there? Theory! Propaganda! How to write papers no
one will read about stuff no one cares about. 100 year old science. At
this point defenders of the System start howling about ``pure
knowledge'', by which they mean theory over ``applied knowledge'', by
which they mean actual knowledge about how the world works. Theory is a
virus, a disease, and a religion, and it has no place here.

What should we learn? Same subjects, but useful. Why can't biology
majors make drugs? Why can't physics majors build a flying drone? Why
can't chemistry majors build a water desalination plant? Many subjects
should, I think, be totally eliminated, as they have no real value, such
as economics and computer ``science''.

\subsection{Science Grad School: The Ponzi
Scheme}\label{science-grad-school-the-ponzi-scheme}

You don't have to go to graduate school to see how much it resembles a
pyramid scam. Suppose any one professor has about five graduate students
at a time, each of whom takes about five years to graduate with a PhD.
If a professor does this for 30 years, they will create on average one
PhD per year or 30 PhD's total. Now suppose all those PhD's find jobs
similar to what their advisor has. This is possible only if the size of
the field increases by 30 times in 30 years. Very crudely this
corresponds to about 12\% per year. Given the expansion of the physical
sciences and their satellites in the years during and directly after
World War II, building a scheme like this in those years actually made
some sense. However, those days are decades behind us, and now as
research budgets shrink and schools, companies, and government agencies
are systematically destroyed by politics, this math looks much more like
a pyramid scam than a response to society's needs.

As a grad student you will \emph{probably} never get the job you're
supposedly being trained for. But you will dedicate 5-7 years of your
life to helping someone who \emph{did} get that job to continue to climb
the academic ladder. The people at the bottom of the academic pyramid
spend their lives working to help the people at the top, and then are
mostly cast aside.

The tenure clock then puts yet another opportunity for exploitation in
the career path, making yet another way for people now well into middle
age to work long hours for more years to build up an academy that they
might then be cast out of.

\subsection{Hollowing Out of the
Academy}\label{hollowing-out-of-the-academy}

There are many factors that have contributed to the downfall of the
university system over the last few years. I would argue that since
Ronald Reagan was elected president of the United States in 1980 there
has been a coordinated ideological war against all culture that might
exist outside of the profit system, and that universities have felt the
brunt of this particularly hard.

A robust, healthy, independent, and publicly supported university system
could provide a real balance against mainstream corporate power if it
existed. It is therefore strategically important for the lords of global
corporate rule that they be as controlled as possible by corporations
and the central government so that they cannot exercise a check on those
forms of power.

\subsection{Intellectual Property}\label{intellectual-property}

Intellectual property deserves its own section here because it has been
so corrosive to the culture of the academy in so many ways. This
manifesto is opposed to all forms of private property, and particularly
intellectual property, but the patenting and copyrighting of work done
in universities is particularly evil.

It is now standard practice for public tax money to be spent to create
knowledge which then goes into papers behind a paywall protected by
brutal copyright enforcement and into patented or even trade secret
knowledge. If the rule of law actually meant anything this would clearly
constitute a criminal theft from the public. The fact that this is
\emph{not} considered a criminal act is in fact a major indication of
the evil nature of the capitalists' so-called ``rule of law''.

\subsection{Potential Paradises}\label{potential-paradises}

How should a university work, after the fall of capitalism and the
takeover of their administration? They must be self-sufficient. They
must grow all their own food and produce all their own energy and
industrial goods. They must be able to maintain the core knowledge
contained in the libraries with equal or superior care to the existing
capitalist system, as that is the heart of the university.

The education of the incoming young people who are now called
``undergraduates'' must be joined with the construction of this new
world. That is, undergraduates will combine learning with labor to
actually physically grow the food, build the robots, build the
industrial production, build the energy infrastructure etc. The graduate
education will combine more advanced and specialized practical work with
more scholarship: writing and reading as it has been done by scholars
for thousands of years. What are now called ``staff'' must be integrated
into this as well, both on the scholarship and on the work and the
development of new methods of work. Bringing the staff into your
insurrection early also creates powerful allies for later on as you
fully take over the school.

\subsection{University Occupations, Phase
2}\label{university-occupations-phase-2}

Any history of any radical movement will inevitably involve student
occupations. Students typically take over some space on campus, keep the
land from the cops, and carry out various protest actions or teach ins
over some number of days or maybe weeks. In some cases they stand down
after actually getting some demands met by the administration. And then
they end, and everyone moves on.

I suggest that as with all occupations one way they can continue is to
create art that can be moved and lasts and is taken from one action to
another, creating a whole narrative that passes along both the ideas and
products of each action to the next, creating a snowball effect. The
purpose of this chapter is to gently nudge people in the direction of
building this kind of infrastructure instead of just running a protest
as long as possible until the inevitable crackdown.

\subsection{Case Study: Your School}\label{case-study-your-school}

\begin{figure}[htbp]
\centering
\includegraphics{images/Sorbonne.png}
\caption{The Sorbonne}
\end{figure}

\begin{figure}[htbp]
\centering
\includegraphics{images/roosevelt.png}
\caption{Map showing location of Roosevelt Island in the East River,
between Manhattan island and the Eastern boroughs of New York City.}
\end{figure}

You have to write this section and integrate it into your version of the
text, if you are connected with a university.

How much land does the school own? Map that, and go walk the land and
figure out what is on it and how it flows. What energy resources are
connected to that land? A creek or river or ocean is your first choice.
In many cases the local river is near but not directly adjoining the
campus. Figure out what a supply chain would look like that would
connect infrastructure you build between the cracks of the property
society on that river with the campus. Can you charge batteries in large
numbers and bring them to campus using courtiers? Can you build hidden
power lines to bring power in? What about water, can you purify that,
bottle it, and bring it in by courier? Study the map you make of your
campus and environs, and figure out what it would look like to get
supply chains built up that will function under martial law, in all out
war with the campus police.

Do not just plan for the short term! If you play chess with the campus
police you will lose; you want to play go. In go, territory is captured
and lost in a dynamic way, and those gains and losses can oscillate
forever in theory. You want to build long term strategy that involves
not just holding an administration building for as long as possible, but
to become a force that flows around campus, rising and falling as you
see opportunities, and fading away. University insurrections have always
suffered from a lack of strategic thought in terms of how land is taken
and held. Realize that as soon as you create any sort of resistance
fixed in geography you are handing advantage to the campus authorities.
You are fighting a guerrilla war, act accordingly.

Also, the long term plan must be partly legal in nature, to exist within
the capitalist state without further military escalation by authorities.
This means you need to recruit from the law school, and have people
studying the corporate governance of your school specifically to figure
out how to carry out a legal coup d'etat in which full legal control of
the university corporation is won by the insurrectionists. University
corporate charters can be very complex and strange and old. Legal
documents like this often have strange legal loopholes which are hard to
find but which can be exploited at the right moment in such a way that
those who believe in the so-called ``rule of law'' can be convinced to
side with the insurrection.

Note that this ``war'' should never be about violence. While the goals
are similar to violent insurrections for a more democratic and free
campus, the methods should never involve competing in violence with
those who have made it there life's work(the campus police and their
backers). Rather, we must win on superior infrastructure. Build better
things: better technology, better classes, better scholarship, better
classrooms, better food, better housing, better libraries. Do that and
you can really win the academy back for the scholars away from the
business scam artists who now run all universities.

\subsection{Universities: a Global
Beachhead}\label{universities-a-global-beachhead}

Taking universities over can be a huge international beach head that
goes around all national borders. Every country has them. And they all
have very energetic young people who want to change the world and try
new things. They can be of huge global importance in the years to come
as the current system falls apart. Capitalism is killing the
universities, and the faculty who love the schools will see eventually
that they have to abandon capitalism if they are to see what is good in
the system continue past this century.
